\documentclass[10pt,a4paper]{article}
\usepackage{amsmath,amssymb,amsthm,graphicx}
\usepackage{hyperref}
\usepackage{enumitem}
\usepackage{mathtools}
\newcommand{\norm}[1]{\left\lVert#1\right\rVert}
\newcommand{\ket}[1]{\left|#1\right\rangle}
\newtheorem{theorem}{Theorem}
\newtheorem*{theorem*}{Theorem}
\newtheorem{definition}[theorem]{Definition}
\newtheorem{corollary}[theorem]{Corollary}
\newtheorem*{corollary*}{Corollary}
\newtheorem{proposition}[theorem]{Proposition}
\newtheorem{lemma}[theorem]{Lemma}
\newtheorem{claim}[theorem]{Claim}
\newtheorem{remark}{Remark}
\newtheorem*{remark*}{Remark}
\title{Exact versus parallel approximate counting equivalence forces a zero-error second-level collapse}
\author{Simon C. Marshall, Scott Aaronson, Vedran Dunjko}
\date{doi:10.4230/LIPIcs.CCC.2025.5; CC BY-SA 4.0}
\begin{document}
\maketitle
\noindent\textit{Editorial note.} Full native author source redistributed under its actual CC BY-SA4.0 license, not relabeled CC BY. The core main theorem native1 corresponds to official2; full author Sections2-3 retained with all support and natural interpretation proof. No mathematical claims corrected. Section4 explicitly declines formal reproduction of the sampling adaptations and is not counted as another complete main result.\par\medskip
\noindent\textit{Editorial note.} Revision r002 inserts only three additional mathematical line breaks in the phi-prime multline in a separate render copy; every original symbol and term is preserved. Native original unchanged.\par\medskip
\noindent\textit{Allowed standard prerequisites (editorial).}\begin{itemize}
\item Random self-reducibility and checkability of counting classes; Toda theorem. Author cited statements of Feigenbaum-Fortnow, Arvind-Kobler and Toda retained verbatim in native Sections2-3; no recursive reproof of those standard prerequisites
\end{itemize}\medskip
\section{Main theorem}\label{sect:Approximate-exact counting theorem}
In this section we prove our main theorem:


\begin{theorem}\label{thm: main}
    If $\mathsf{P^{\# P}=BPP_{\parallel}^{NP}}$ then $\mathsf{P^{\# P}=ZPP^{NP}}$.
\end{theorem}



We will assume readers are broadly familiar with classes such as $\mathsf{\# P}$, where classes are undefined we use the definitions given in the complexity zoo \cite{zoo}. We use the notation $\mathsf{A_{\parallel}^{B}}$ for $\mathsf{A}$ with one set of parallel oracle calls to $\mathsf{B}$.

It is useful to first give an informal description of the proof of theorem \ref{thm: main} before we proceed to the formal proof to give the reader a better idea of the purpose of each lemma. Broadly the proof relies on the notion of random-reducibility, where a given instance of a problem can be solved by a polynomial time algorithm with randomly selected instances of some other problem. If a problem can be randomly reduced to itself we say it is random self-reducible. 
The first key realisation for Theorem \ref{thm: main} is that a random self-reducible language that is in $\mathsf{BPP^{NP}_{\parallel}}$ is itself random reducible to $\mathsf{NP}$, because randomly selecting $x$ for a language in $\mathsf{BPP^{NP}_{\parallel}}$ gives rise to random set of non-adaptive $\mathsf{NP}$ calls, which meets the definition of random reducibility to $\mathsf{NP}$.
Next, as $\mathsf{\#P}$ is random self-reducible \cite{feigenbaum1990random}, the antecedent of Theorem \ref{thm: main} implies that it is random reducible to $\mathsf{NP}$.
The second key realisation is to notice that the proof by Feigenbaum and Fortnow that $\mathsf{NP}$ cannot be random self-reducible unless it is also in $\mathsf{coNP/poly}$ extends beyond self reducibility; any language that is random reducible to $\mathsf{NP}$ is in $\mathsf{coNP/poly}$. 
By repeating the argument for $\mathsf{coNP}$ we show $\mathsf{\#P\subseteq NP/poly \cap coNP/poly}$. Arvind and Köbler \cite{ARVIND2002257} showed any checkable language in $\mathsf{NP/poly \cap coNP/poly}$  is low for $\mathsf{\Sigma^P_2}$ which gives us a polynomial hierarchy collapse to $\mathsf{\Sigma_2^P}$. The final collapse to $\mathsf{ZPP^{NP}}$ comes from using the $\mathsf{BPP^{NP}_{\parallel}}$ algorithm to produce proofs verifiable in $\mathsf{P^{NP}}$, thereby achieving zero error.

We begin by defining random reducibility. A problem is randomly self-reducible if any given instance of the problem can be solved probabilistically by a polynomial time algorithm with access to solved random instances of the same problem.  
For our purposes we will work with a similar concept: random-reducibility, which is the same except the random instances may be from a different problem. It should be stated that this is different from a `randomized reduction', which is a randomised Karp reduction. Random reductions are sometimes called 1-locally random reductions to avoid this confusion.

\begin{definition}[Random Reducible]\label{def: random reducible}
    A function $f$ is \textbf{nonadaptively $k(n)$ random-reducible} to a function $g$ if there are polynomial time computable functions $\sigma$ and $\phi$ with the following two properties.
    \begin{enumerate}[label={(\arabic*)}]
        \item For all $n$ and all $x \in \{0,1\}^n$ $f(x)$ can probably be solved by $\phi$ with instance of $g(y)$ for $y$ selected by $\sigma$:
            \[
            \Pr_{r}\big(
            f(x)=\phi(x,r,g(\sigma(1,x,r)),\ldots, g(\sigma(k, x, r)))
            \big)\geq2/3
            \]
        \item For all $n$, all pairs $\{x_1,x_2\}\subset \{0,1\}^n$ and all $i$, if $r$ is chosen uniformly at random then $\sigma(i,x_1,r)$ and $\sigma(i,x_2,r)$ are identically distrubuted.
    \end{enumerate}
    If both conditions hold we say ($\sigma$, $\phi$) is a random-reduction from $f$ to $g$.
\end{definition}

As we are dealing with no other forms of random reducibility we will just say `$f$ is rr to $g$' when $f$ is non-adaptively $k(n)$ random-reducible to $g$ for some polynomial $k$. If a function is rr to itself then we say the function is random self-reducible (rsr). A set is rr to another if its characteristic function is rr. A set of languages, $\mathsf{A}$ is rr to $\mathsf{B}$ if every $L\in \mathsf{A}$ is rr to some $L'\in\mathsf{B}$.

As with probabilistic classes like $\mathsf{BPP}$, we can boost random reducibility an arbitrarily low ($2^{-n}$) failure probability.

\begin{lemma}[\cite{feigenbaum1990random}]\label{lma: boost rr}
    If function $f$ is non-adaptively $k(n)$ random-reducible to $g$ then for all $t(n)$, $f$ is non-adaptively $24t(n)k(n)$ random-reducible to $g$ where condition (1) holds for at least $1-2^{-t(n)}$ of the r's in $\{0,1\}^{24t(n)k(n)}$.
\end{lemma}

\begin{remark*}
    In general, this boosting may not work for a definition of random reducibility dealing with relations instead of functions as they may have multiple correct outputs. However, counting problems like $\mathsf{Perm}$, which has only one correct output for a given input, can be boosted.
\end{remark*}


\begin{lemma}[\cite{feigenbaum1990random}]\label{lma: perm is rsr}
    Any $\mathsf{\#P}$-complete language is rsr.
\end{lemma}

This concludes our definitions, we can now state an intermediate theorem to our final result.

\begin{theorem}\label{thm: BPP is rr to NP}
    If $\mathsf{P^{\# P}= BPP_{\parallel}^{NP}}$ then $\mathsf{P^{\# P}}$ is random reducible to $\mathsf{NP}$.
\end{theorem}

Before we prove this, we will provide a number of smaller results. In the following lemma the notation $\mathsf{A^{(B[1])}}$ means $\mathsf{A}$ with one oracle call to $\mathsf{B}$.

\begin{lemma}\label{lma: P to 1 sharp P is rr to 1 sharp P}
    Any language in $\mathsf{P^{\# P[1]}}$ is random reducible to $\mathsf{\# P}$.
\end{lemma}
\begin{proof}
    As $\mathsf{\# P}$-complete languages are random self-reducible and only one oracle call is needed we can use a random reduction of $\mathsf{\# P}$ to give the answer to the oracle call and use the polynomial time $\phi$ algorithm to perform the rest of the $\mathsf{P}$ algorithm.
\end{proof}

The next lemma captures the intuition that $\mathsf{BPP^{NP}_{\parallel}}$ algorithms are `almost' rr to $\mathsf{NP}$, only missing the random element selection part of the definition (property (2)).

\begin{lemma}\label{lma:BPP rr similarity}
    For all $L$ in $\mathsf{BPP^{NP}_{\parallel}}$ there exists polynomial time computable functions $\sigma_{B}$, $\phi_{B}$ such that
    \[
    \Pr_{r}\left(
            L(x)=\phi_{B}(x,r,\mathsf{SAT}(\sigma_{B}(1,x,r)),\ldots, \mathsf{SAT}(\sigma_{B}(m, x, r)))
            \right)\geq 1-2^{-n}.
    \]
\end{lemma}
\begin{proof}
    If $L\in \mathsf{BPP^{NP}_{\parallel}}$ then there exists a polynomial time algorithm, $A$, to decide $x \in L$ which calls only one set of up to $m$ non-adaptive $\mathsf{NP}$ queries. We assume the $\mathsf{NP}$ calls $\mathsf{SAT}$ for simplicity. Let the algorithm for $\sigma_{B}(i,x,r)$ run $A$ until the step involving oracle queries and then output just the $i$'th query (assuming some arbitrary query ordering).
    The algorithm for $\phi_{B}$ is now just the algorithm $A$ using the oracle calls produced asked by $\sigma_{B}(i,x,r)$ in place of directly making its own calls. Since both $\phi_{B}$ and $\mathsf{\sigma_{B}}$ have access to the same random string $r$ they will both select the same oracle calls.
\end{proof}
Lemma \ref{lma:BPP rr similarity} is useful as it proves that $\mathsf{BPP^{NP}_{\parallel}}$ fulfils property (1) of the definition of random reducibility to $\mathsf{NP}$, but does not prove the random distribution of ${\sigma_B(1,x,r)}$. The next theorem shows that if a langauge already randomly reduces to $\mathsf{BPP^{NP}_{\parallel}}$ then this random reduction condition (condition (2)) is also fufilled, and thus the language is random reducible to $\mathsf{NP}$

\begin{lemma}\label{lma: rr to BPP NP parallel is rr to NP}
    All languages random reducible to a language in $\mathsf{BPP^{NP}_{\parallel}}$ are random reducible to $\mathsf{NP}$.
\end{lemma}
\begin{proof}Let $L$ be a language which is random reducible to $L_B\in \mathsf{BPP^{NP}_{\parallel}}$. We will demonstrate that $L$ is rr to $\mathsf{NP}$ by writing out the definition of random reducibility to ${L_B}$ then using lemma \ref{lma:BPP rr similarity} to substitute the calls to ${L_B}$ with a formula using calls to $\mathsf{SAT}$. We can then rearange this into a format which directly makes calls to $\mathsf{SAT}$ and we show that these calls inherit the randomness property from the random reduction of $L$ to $L_B$.

    Let us assume the random reduction of $L$ to $L_B$ on input $x$ uses $k$ calls, from the definition of random reducibility there exists $\phi$ and $\sigma$ such that
    \[
\Pr_{r}\left(
            L(x)=\phi(x,r,L_{\mathsf{BPP}}(\sigma(1,x,r)),\ldots, L_{{B}}(\sigma(k, x, r)))
            \right)\geq 1-2^{-n}.
\]

Define $y_i:=\sigma(i, x, r)$. 

\[
\Pr_{r}\left(
            L(x)=\phi(x,r,L_{\mathsf{B}}(y_1),\ldots, L_{\mathsf{BPP}}(y_k))
            \right)\geq 1-2^{-n}
\]

Using lemma \ref{lma:BPP rr similarity} we can substitute $L_B$ with $\phi_B$, $\sigma_B$ and a random string $r_i$ for the i'th call to $L_B$. In the following we assume each lemma-\ref{lma:BPP rr similarity}-reduction uses $m$ $\mathsf{SAT}$ calls.

\begin{align*}
\Pr_{r, r_1, \ldots r_k}\big(
            L(x)=\phi(x,r,&\phi_{B}(y_1,r_1,\mathsf{SAT}(\sigma_{B}(1,y_1,r_1)),
            \ldots, \mathsf{SAT}(\sigma_{B}(m, y_1, r_1))),\\ \ldots, &\phi_{B}(y_k,r_k,\mathsf{SAT}(\sigma_{B}(1,y_k,r_k)), \ldots, \mathsf{SAT}(\sigma_{B}(m,y_k,r_k)))
            \big)
            \geq 1-(k+1)2^{-n}.
\end{align*}

The failure probability $1-(k+1)2^{-n}$ comes the failure probability of the $k$ reductions combined with the failure probability from the random reduction.

We can now begin to combine elements of the reduction from $L$ to $L_B$ with reduction from $L_B$ to $\mathsf{SAT}$. We start with $\phi$ which we combine with $\phi_B$ to define $\phi'$, informally we can think of this as doing the two polynomial-time algorithms as a larger polynomial time algorithm.
\begin{multline*}
\phi'(x,r,r_1,\ldots,r_k,\\
\mathsf{SAT(\sigma_B(1,y_1,r_1})),\ldots \mathsf{SAT(\sigma_B(m,y_k,r_k}))):=
\\
\phi(x,r,\phi_{B}(y_1,r_1,\mathsf{SAT}(\sigma_{B}(1,y_1,r_1)),
\ldots, \mathsf{SAT}(\sigma_{B}(m, y_1, r_1))),\ldots,\\
\phi_{B}(y_k,r_k,\mathsf{SAT}(\sigma_{B}(1,y_k,r_k)),\\
\ldots, \mathsf{SAT}(\sigma_{B}(m,y_k,r_k))
\end{multline*}

We use `$x\#y$' to denote $y$ appended to $x$. Define $\mathbf{r}:=r\#r_1\#\ldots\#r_k$.

Again we define a new function, $\sigma'$, as the combination of two existing function, $\sigma$ and $\sigma_B$:

\[
\sigma'(i,j,x,\mathbf{r}):= \sigma_B(i,\sigma(j,x,r),r_j).
\]

Which gives the following simplified equation:
\begin{equation}\label{eq: finally}
\Pr_{\mathbf{r}}\big(
            L(x)=
\phi'(x,\mathbf{r}, \mathsf{SAT}(\sigma'(1,1,x,\mathbf{r})),\ldots \mathsf{SAT}(\sigma'(m,k,x,\mathbf{r}))
            \big)
            \geq 1-(k+1)2^{-n}.
\end{equation}

We will now directly check the definition of $\sigma'$, $\phi'$ and the equation \ref{eq: finally} against the definition of random-reducbile, recall the two conditions definition \ref{def: random reducible}:
\begin{enumerate}[label={(\arabic*)}]
        \item For all $n$ and all $x \in \{0,1\}^n$ $f(x)$ can probably be solved by $\phi$ with instance of $g(y)$ for $y$ selected by $\sigma$:
            \[
            \Pr_{r}\big(
            f(x)=\phi(x,r,g(\sigma(1,x,r)),\ldots, g(\sigma(k, x, r)))
            \big)\geq2/3
            \]
        \item For all $n$, all pairs $\{x_1,x_2\}\subset \{0,1\}^n$ and all $i$, if $r$ is chosen uniformly at random then $\sigma(i,x_1,r)$ and $\sigma(i,x_2,r)$ are identically distrubuted.
    \end{enumerate}

We can see that these two conditions are met: 

\begin{enumerate}[label={(\arabic*)}]
        \item For sufficiently large $n$, $1-(k+1)2^{-n}>2/3$ therefore equation \ref{eq: finally} is of the form:
            \[
            \Pr_{r}\big(
            f(x)=\phi(x,r,\mathsf{SAT}(\sigma(1,x,r)),\ldots, \mathsf{SAT}(\sigma(k, x, r)))
            \big)\geq2/3.
            \]
        \item For all $n$, all pairs $\{x_1,x_2\}\subset \{0,1\}^n$, all $i$ and $j$, if $r$ is chosen uniformly at random then $\sigma(j,x_1,r)$ and $\sigma(j,x_2,r)$ are identically distributed (as per their definition), therefore $\sigma'(i\#j,x_{1/2},\mathbf{r})$ is identically distributed too. This is because the only dependence on $x$ is through the identically distributed $\sigma$: $\sigma'(i\#j,x,\mathbf{r})= \sigma_\mathsf{B}(i, \sigma(j, x, r), r_j)$. This holds for all $r_i$, $r_j$. 
    \end{enumerate}
Thus fufilling the conditions for $L$ to be random reducible to $\mathsf{NP}$
\end{proof}

Finally, we can proceed with the proof of theorem \ref{thm: BPP is rr to NP}

\begin{proof}[Proof of theorem \ref{thm: BPP is rr to NP}]
Toda \cite{toda1991pp} showed that $\mathsf{PH\subseteq P^{\# P[1]}}$. Therefore, by the assumption of the theorem if $\mathsf{P^{\# P}= BPP_{\parallel}^{NP}=PH}$ then $\mathsf{P^{\# P} = P^{\# P[1]}}$. 

By lemma \ref{lma: P to 1 sharp P is rr to 1 sharp P} we know $\mathsf{P^{\# P}}$ is rr to $\mathsf{\# P}$. Under the assumption of this theorem, $\mathsf{\# P \subseteq BPP^{NP}_{\parallel}}$, $\mathsf{P^{\# P}}$ is rr to $\mathsf{BPP^{NP}_{\parallel}}$. By lemma \ref{lma: rr to BPP NP parallel is rr to NP} this means $\mathsf{P^{\# P}}$ is rr to $\mathsf{NP}$.
\end{proof}

The next result is heavily based on the work of Feigenbaum and Fortnow \cite{feigenbaum1990random}. They show that if $\mathsf{NP}$ is random self reducible then $\mathsf{coNP}$ is in $\mathsf{NP/poly}$, while the result we are trying to show is slightly different than this, the method is very similar.

\begin{theorem}\label{thm: rr to NP is NP/poly}
    Any language that is random-reducible to a language in $\mathsf{NP}$ is also in $\mathsf{coNP/poly \cap NP/poly}$.
\end{theorem}

\begin{proof}
    Let $\mathsf{AM^{poly}}$ be $\mathsf{AM}$ with polynomial advice given to Arthur\footnote{ This is not the same as $\mathsf{AM/poly}$, as the latter must be an $\mathsf{AM}$ language for all advice, whereas the former only needs to be defined for the correct advice.}. It turns out that $\mathsf{AM^{poly}=NP/poly}$ and $\mathsf{coAM^{poly}=coNP/poly}$ \cite{feigenbaum1990random}.
    We will prove $L$ is in $\mathsf{AM^{poly}\cap coAM^{poly}}$, beginning with $L \in \mathsf{AM^{poly}}$.

    Suppose $L(x)$ is non-adaptively $k(n)$ random-reducible to $\mathsf{SAT}(x)$ with error probability $2^{-n}$. We will adopt the notation $y_i=\sigma(i,x,r)$. For instance size $n$ we will give Arthur the advice $(p_1,\ldots, p_k)$ where $p_i$ is the probability that $y_i=1$,
    \[
    p_i=\Pr_{r}(\mathsf{SAT}(\sigma(i, x,r))=1).
    \]
    As $\sigma(i, x,r)$ is distributed identically (given uniform random $r$) for all $x$ this advice does not depend on the input $x$.

    The proof system will consist of Arthur selecting $m=9k^3$ random strings, $\{r_j\}_{j\in[m]}$, and passing this to Merlin, these random strings defines which $\mathsf{SAT}$ queries, $y_{0,j},\ldots, y_{k,j}$, will be made. Merlin will hand back $m$ `transcripts'. These transcripts consist of a list $w_{i,j}$ which is either a witness for $y_{i,j}$ or is the string `NIL' (which is interpreted as the claim that $y_{i,j}\notin \mathsf{SAT}$). For each of the $m$ random strings we get the transcript:
    \[
    Transcript(x,r_j)= (w_{0,j},w_{1,j}, \ldots, w_{k,j})
    \]
    Arthur performs three checks:
    \begin{enumerate}[label={(\arabic*)}]
        \item For all $i,j$ either the witness $w_{i,j}$ is `NIL' or he checks the witness is valid for $y_{i,j}\in \mathsf{SAT}$
        \item Let $b_{i,j}$ be 0 if $w_{i,j}$ is `NIL' and one otherwise. For all $j\in [m]$ Arthur checks that $\phi(x,r_j, b_{0,j}, \ldots, b_{k,j})=1$, i.e. If Merlin has told the truth then $\phi$ will accept.
        \item For each $i\in [k]$ Arthur checks that more than $p_im-2\sqrt{km}$ of the $y_{i,j}$ have been proved to be in $\mathsf{SAT}$, if this condition does not hold he rejects.
    \end{enumerate}
    If these checks all pass, then Arthur accepts. 
    The trick in this proof is this final condition. Over the random strings, each $y_i$ has some probability of being in $\mathsf{SAT}$ with a given variance, by picking $m$ to be large we force this variance to be small. With high probability, a correct answer lies in this range. Since Merlin cannot lie about yes instances (since he must prove them), he can only lie about no instances. However, if he lied too much it would be clear that not enough of $y_i$ are in $\mathsf{SAT}$, thus we could probabilistically guess he was cheating. To exploit this, $m$ is picked precisely so Merlin cannot cheat on all $j$ without exceeding his `lying budget'. We will now formally show soundness and completeness.

    \textit{Completness.} If $x\in L$ and Merlin is honest, providing witnesses for all $y_{i,j}$ in $\mathsf{SAT}$, condition (1) will always hold. Condition (2) holds with probability at least $1-2^{-n}$ (by the definition of rr). Therefore we just need to show condition (3) holds with high probability.

    Let $Z_{i,j}:=(y_{i,j}=1)$ and $Z_{i}=\sum_{j=1}^m Z_{i,j}$. We defined our advice so $p_i= \mathbb{E}[Z_i]/m$, we can derive that $\operatorname{Var}(X)=p_i(1-p_i)m<m$. We can use these properties and Chebyshev's inequality to bound our probability of failing check (2).
    \begin{align*}
        \Pr(Z_i\leq p_im-2\sqrt{km})\\
        \leq\Pr(\norm{Z_i-p_im}\geq 2\sqrt{km})\\
        =\Pr(\norm{Z_i-\Bar{Z_i}}\geq 2\sqrt{km})\\
        \leq \operatorname{Var}(X)/4km\leq 1/(4k)\leq1/4
    \end{align*}
    The probability of failing any check given $x\in L$ is less than $\frac{1}{4}+2^{-n}$ which is less than $1/3$ for large enough $n$. This proves completeness

    \textit{Soundness.} Suppose $x \notin L$, if Arthur accepts then all 3 checks must have passed for all $j\in[m]$. If some $y_{i,j}$ is in $\mathsf{SAT}$ but Merlin has provided $w_{i,j}=`NIL'$ then we say Merlin has lied about $y_{i,j}$, if check (1) passes Merlin has not lied about any `yes' instances, so we disregard this case.
    The probability of the random reduction failing is $2^{-n}$ without any lies, so for all $m$ reductions to fail Merlin must lie $m$ times with probability $(1-2^{-n})^m> 1-m2^{-n}$.

    If Merlin lies $m$ times there must be an $i$ that he claims at least $m/k$ of the $y_{i,j}$ are not in $\mathsf{SAT}$ when they are. To satisfy condition (3) at least $p_im-2\sqrt{km}\mathbf{+m/k}$ of the $y_{i,j}$ must be in $\mathsf{SAT}$ to leave `room' for Merlin to lie $m/k$ times without violating (3). The probability of this is given by a Chernoff bound on the random variable $Z_i$, as defined above.
    \begin{align*}
    \Pr\left(Z_i\geq p_im-2\sqrt{km}+m/k\right)\\
    =\Pr\left(Z_i-p_im\geq m/k-2\sqrt{km}\right)\\
    =\Pr\left(Z_i-p_im\geq 9k^3/k-2\sqrt{k9k^3}\right)\\
    =\Pr\left(Z_i-p_im\geq 3k^2\right)\\
    \leq e^{-2\times 9k^4/9k^3}=e^{-2\times k}\\
    \leq 1/(4k)
    \end{align*}
    Combining this with the normal probability that the random reduction fails even with correct oracle answers (which had a $2^{-n}$ probability) gives an acceptance probability given $x\notin L$ of less than $m2^{-n}+1/(4k)<1/3$ for large enough $n$. This proves soundness.
    
    The previous analysis can just as easily be repeated for $\mathsf{coAM^{poly}}$ with Merlin proving no instances, giving $L\in \mathsf{AM^{poly} \cap coAM^{poly}}$. By Feigenbaum and Fortnow we know this equals $\mathsf{NP/poly \cap coNP/poly}$.
\end{proof}

To prove the next Theorem we must recall a crucial result by Arvind and Köbler, which checkability \cite{ARVIND2002257,zoo}. The notion of checkability is somewhat involved but intuitively has to do with whether efficent programs that decide the set can themselves be efficently checked. Since we will just need the checkability of $\mathsf{PP}$ and $\mathsf{\# P}$ to deploy the following theorem we will not formally define it, instead we refer the reader to \cite{ARVIND2002257, blum1995designing}.

\begin{theorem*}[\cite{ARVIND2002257}, Theorem 22]
    Checkable sets in $\mathsf{NP/poly \cap coNP/poly}$ are low for $\mathsf{\Sigma_2^P}$.
\end{theorem*}

This result is imperative to proceed, as while it is true that $\mathsf{NP\subset (NP\cap coNP)/poly}$ implies $\mathsf{PH=ZPP^{NP}}$, the same is not true for $\mathsf{NP\subset (NP/poly)\cap (coNP/poly)}$.

\begin{lemma}\label{lma: sharp P in conp/poly}
    If $\mathsf{P^{\# P}}\subset (\mathsf{coNP/poly \cap NP/poly})$ then $\mathsf{P^{\# P}=\Sigma_2^P\cap \Pi_2^P}$
\end{lemma}
\begin{proof}
By Toda $\mathsf{P^{\#P}=P^{PP}}$ \cite{toda1991pp}.    
    If $\mathsf{P^{PP}}$ is in $\mathsf{\mathsf{coNP/poly \cap NP/poly}}$ then clearly so is $\mathsf{PP}$, combining this with the existence of $\mathsf{PP}$-complete checkable sets \cite{ARVIND2002257} we know that $\mathsf{PP}$ is low for $\mathsf{\Sigma_2^P}$.

    This lowness implies we can put  $\mathsf{P^{\# P}}$ in $\mathsf{\Sigma_2^P}$ by the following series of inclusions
    \[
    \mathsf{
    P^{\#P}=P^{PP}\subseteq (\Sigma_2^P)^{PP} = \Sigma_2^P.
    }
    \]

    As $\mathsf{\Sigma_2^P\subseteq P^{\# P}}$ by Toda \cite{toda1991pp} we can fix the previous inclusion into an equality: $\mathsf{\Sigma_2^P =  P^{\# P}}$.

    As $\mathsf{\Pi_2^P \subseteq P^{\# P}}$ (Toda \cite{toda1991pp}) we get $\mathsf{\Pi_2^P\subseteq \Sigma_2^P}$.
    Which gives us the final equality:
    \[
    \mathsf{P^{\# P} \subseteq \Sigma_2^P \subseteq P^{\Pi_2^P}= P^{\Sigma_2^P\cap\Pi_2^P}=\Sigma_2^P\cap\Pi_2^P}
    \]
\end{proof}

The previous lemma has achieved a collapse to the second level but we can collapse to $\mathsf{ZPP^{NP}}$ by using the proveability of languages in $\mathsf{\Sigma_2^P\cap \Pi_2^P}$ to make the $\mathsf{BPP^{NP}_{\parallel}}$ algorithm zero-error.

\begin{lemma}\label{lma: sharp P in ZPP}
    All languages in $\mathsf{BPP^{NP}\cap\Sigma_2^P\cap \Pi_2^P}$ are in $\mathsf{ZPP^{NP}}$
\end{lemma}
\begin{proof}
    Fix some $L\in \mathsf{BPP^{NP}\cap\Sigma_2^P\cap \Pi_2^P}$. We can check if $x\in L$ using the $\mathsf{BPP^{NP}}$ algorithm. Use the $\mathsf{BPP^{NP}}$ program to determine a proof of this fact for either the $\mathsf{\Sigma_2^P}$ verifier or the $\mathsf{\Pi_2^P}$ verifier. We can test either proof in $\mathsf{P^{NP}}$ using the verifiers from either $\mathsf{\Sigma_2^{P}}$ or $\mathsf{\Pi_2^{P}}$. With probability $1-2^{-n}$ the $\mathsf{BPP^{NP}}$ algorithm succeeds in producing a valid proof in polynomial time, any valid proof proves or disproves $x\in L$. Therefore in polynomial time the algorithm has successfully determined if $x\in L$ with probability $1-2^{-n}$ and never incorrectly answers, placing $L\in \mathsf{ZPP^{NP}}$.
\end{proof}

This completes the proof of our main theorem which is given below.

\begin{theorem*}
    If $\mathsf{P^{\# P} = BPP^{NP}_{\parallel}}$ then $\mathsf{P^{\# P} = ZPP^{NP}}$
\end{theorem*}
\begin{proof}[Summary of proof]
    If $\mathsf{P^{\# P} = BPP^{NP}_{\parallel}}$ then $\mathsf{P^{\#P}}$ is random reducible to $\mathsf{NP}$ (Theorem \ref{lma: perm is rsr}).

    If $\mathsf{P^{\# P}}$ is random reducible to $\mathsf{NP}$ then $\mathsf{P^{\#P}}$ is in $\mathsf{coNP/poly \cap NP/poly}$ (Theorem \ref{thm: rr to NP is NP/poly}).

    If $\mathsf{P^{\#P}\subseteq coNP/poly \cap NP/poly}$ then $\mathsf{P^{\#P}} = \mathsf{\Sigma_2^P\cap \Pi_2^P}$ (Lemma \ref{lma: sharp P in conp/poly}).

    If $\mathsf{P^{\#P}} = \mathsf{\Sigma_2^P\cap \Pi_2^P}$ and $\mathsf{P^{\#P}= BPP^{NP}}$ then $\mathsf{P^{\#P}= ZPP^{NP}}$ (Lemma \ref{lma: sharp P in ZPP}).
\end{proof}

\begin{remark*}
    Interestingly we can perform all of the oracle calls required to produce the proof of $x\in L$ in one parallel step, it then only takes one extra Oracle call to check this proof. Thus it is possible to answer correctly with probability $1-2^{-n}$ or with `I don't know' after only 2 rounds of parallel $\mathsf{NP}$ oracle calls.
\end{remark*}

\begin{remark*}
    Feigenbaum and Fortnow \cite{feigenbaum1990random} showed their proof applies beyond non-adaptive random self-reducibility, up to logarithmically many adaptive steps are allowed. For the same reasons, our proof could be extended to logarithmically many rounds of polynomially many parallel oracle calls.
\end{remark*}

The following theorem formalises the `natural interpretation' of our result given in the intro.

\begin{theorem}[Natural interpretation]\label{thm: natural}
    If $\mathsf{P^{ExactCount}_{\parallel}}=\mathsf{P^{ApproxCount}_{\parallel}}$ then $\mathsf{P^{\#P}=ZPP^{NP}}$
\end{theorem}

This result becomes easy to show after Theorem \ref{thm: postbqp postbpp}, so we will prove it at the end of the next section.



%%%%%%%%%%%%%%%%%%%%%%%%%%
%%%%%%%%  postBQP %%%%%%%%
%%%%%%%%%%%%%%%%%%%%%%%%%%
\section{PostBQP and PostBPP}\label{sect: postBQP v postBPP}

An immediate consequence of the last section's main theorem is to the question of $\mathsf{PostBQP\stackrel{?}{=} PostBPP}$, i.e. is a poly-time quantum computer with postselection equal to a classical computer with postselection? The best existing answer was given when Aaronson showed $\mathsf{PostBQP=PP}$ \cite{aaronson2014equivalence}, thus equality would imply $\mathsf{PH=BPP^{NP}}$ (a collapse to the third level). However, as $\mathsf{PostBPP}\subseteq \mathsf{BPP^{NP}_{\parallel}}$ \cite{o2018weakness, han1997threshold} our theorem can be used to improve this collapse to the second level. Before we provide proof of this we formally define $\mathsf{PostBQP}$ and the operator $\mathsf{BP\cdot}$ which makes a complexity class probabilistic.

\begin{definition}
$\mathsf{PostBQP}$ is the set of languages for which there exists a uniform family of polynomial width and polynomial depth quantum circuits $\{C_n\}_ {n\geq 1}$ such that for any input x,
\begin{enumerate}[label=(\roman*)]
    \item There is a non-zero probability of measuring the first qubit of $C_n \ket{0\ldots 0}\ket{x}$ to be $\ket{1}$.
    \item If $x\in L$, conditioned on the first qubit being $\ket{1}$, the second qubit is $\ket{1}$ with probability at least 2/3.
    \item If $x\notin L$, conditioned on the first qubit being $\ket{1}$, the second qubit is $\ket{1}$ with probability at most 1/3.
\end{enumerate}
\end{definition}

$\mathsf{PostBPP}$ can be defined similarly with classical circuits and additional input of random bits.

\begin{definition}
    Let K be a complexity class. A language $L$, is in $\mathsf{BP\cdot K}$ if there exists a language $\mathsf{A \in K}$ and a polynomial $p$ such that for all strings $x$.
$$Pr(r \in \{0,1\}^{p(|x|)}: x \in L \iff x\#r \in A) \geq 2/3$$
\end{definition}

This covers the necessary definitions and we can proceed with the collapse result of interest.

\begin{theorem}\label{thm: postbqp postbpp}
    If $\mathsf{PostBQP=PostBPP}$ then $\mathsf{P^{\#P}=ZPP^{NP}}$ and the polynomial hierarchy collapses to the second level.
\end{theorem}

\begin{proof}
Assume $\mathsf{PostBQP=PostBPP}$. By Toda $\mathsf{PH\subseteq P^{\#P}}$, by Aaronson \cite{aaronson2005postBQP} $\mathsf{PostBQP=PP}$ and by $\mathsf{PostBPP \subseteq BPP^{NP}_{\parallel}}$ \cite{o2018weakness, han1997threshold}. This gives: 
\[
\mathsf{
PP \subseteq BPP^{NP}_{\parallel} \text{  and  }  P^{\#P} = P^{PP} \subseteq P^{BPP^{NP}} = BPP^{NP} = PH \subseteq P^{\#P}.
}
\]

Therefore $\mathsf{P^{\#P}=PH}$.

At this point we recall an interesting lemma from Toda \cite{toda1991pp} to continue.
\[\mathsf{
PP^{PH}\subseteq BP\cdot PP
}\]

Clearly $\mathsf{PH\subseteq PP^{PH}}$ and it can be shown that if $\mathsf{PP\subseteq BPP^{NP}_{\parallel}}$ then $\mathsf{BP \cdot PP \subseteq BPP^{NP}_{\parallel}}$ (a full proof follows in lemma \ref{lma: BP PP is in parallel too}), therefore

\[\mathsf{
PH\subseteq BP\cdot PP \subseteq BPP^{NP}_{\parallel}
}.\]

Since $\mathsf{PH = P^{\#P}}$ and $\mathsf{BPP^{NP}_{\parallel}\subseteq PH}$ we can conclude $\mathsf{P^{\#P}=BPP^{NP}_{\parallel}}$. We now have the condition $\mathsf{P^{\#P}= BPP^{NP}_{\parallel}}$ so by theorem \ref{thm: main} we get $\mathsf{P^{\#P} = ZPP^{NP}}$, the final result.
\end{proof}


\begin{lemma}\label{lma: BP PP is in parallel too}
    If $\mathsf{PP}$ is in $\mathsf{BPP^{NP}_{\parallel}}$ then so is $\mathsf{BP \cdot PP}$
\end{lemma}
\begin{proof}
Fix some $L\in \mathsf{BP\cdot PP}$. By the definition, there exists $A\in\mathsf{PP}$ which defines $L$.
Assume $\mathsf{PP}$ is in $\mathsf{BPP^{NP}_{\parallel}}$, therefore $A$ is in $\mathsf{BPP^{NP}_{\parallel}}$. Therefore $L$ is in $\mathsf{BP\cdot BPP^{NP}_{\parallel}}$.

Note that majority-vote amplification can be used on $\mathsf{BP\cdot PP}$ and $\mathsf{BPP^{NP}_{\parallel}}$ to decrease the probability of failure to less than $1/6$.

If we combine the random coin flips of both the $\mathsf{BP \cdot PP}$ and $\mathsf{BPP}$ parts of the algorithm we get a single failure probability below 1/3, which fits the definition of $\mathsf{BPP^{NP}_{\parallel}}$.
\end{proof}

Theorem \ref{thm: postbqp postbpp} provides a quick a simple proof of Theorem \ref{thm: natural} (the natural interpretation of our result).

\begin{proof}[Proof of Theorem \ref{thm: natural}]
By definition $\mathsf{\#P=ExactCount}$. O'Donnell and Say show $\mathsf{P^{ApproxCount}_{\parallel}=PostBPP}$ \cite{o2018weakness}. 
Therefore we can rewrite the antecedent of the Theorem as $\mathsf{P^{\#P}_\parallel=PostBPP}$.

Any problem in $\mathsf{PP}$ can be solved with one $\mathsf{\#P}$ query, which can obviously be done in one parallel round of oracle calls, so $\mathsf{PP\subseteq P^{\#P}_\parallel}$. This gives $\mathsf{PP=PostBPP}$ (as we already know $\mathsf{PostBPP\subseteq PP}$ and by Theorem \ref{thm: postbqp postbpp} $\mathsf{P^{\#P}=ZPP^{NP}}$.
\end{proof}

\begin{thebibliography}{10}

\bibitem{aaronson2005postBQP}
Scott Aaronson.
\newblock Quantum computing, postselection, and probabilistic polynomial-time.
\newblock {\em Proceedings of the Royal Society A: Mathematical, Physical and Engineering Sciences}, 461(2063):3473--3482, 2005.

\bibitem{aaronson2014equivalence}
Scott Aaronson.
\newblock The equivalence of sampling and searching.
\newblock {\em Theory of Computing Systems}, 55(2):281--298, 2014.

\bibitem{aaronson2011computational}
Scott Aaronson and Alex Arkhipov.
\newblock The computational complexity of linear optics.
\newblock In {\em Proceedings of the forty-third annual ACM symposium on Theory of computing}, pages 333--342, 2011.

\bibitem{zoo}
Scott Aaronson, Greg Kuperberg, and Christopher Granade.
\newblock The complexity zoo, 2005.

\bibitem{aharonov2023polynomial}
Dorit Aharonov, Xun Gao, Zeph Landau, Yunchao Liu, and Umesh Vazirani.
\newblock A polynomial-time classical algorithm for noisy random circuit sampling.
\newblock In {\em Proceedings of the 55th Annual ACM Symposium on Theory of Computing}, pages 945--957, 2023.

\bibitem{ARVIND2002257}
V.~Arvind and Johannes Köbler.
\newblock New lowness results for zppnp and other complexity classes.
\newblock {\em Journal of Computer and System Sciences}, 65(2):257--277, 2002.

\bibitem{blum1995designing}
Manuel Blum and Sampath Kannan.
\newblock Designing programs that check their work.
\newblock {\em Journal of the ACM (JACM)}, 42(1):269--291, 1995.

\bibitem{bremner2016average}
Michael~J Bremner, Ashley Montanaro, and Dan~J Shepherd.
\newblock Average-case complexity versus approximate simulation of commuting quantum computations.
\newblock {\em Physical review letters}, 117(8):080501, 2016.

\bibitem{feigenbaum1990random}
J~Feigenbaum and L~Fortnow.
\newblock On the random-self-reducibility of complete sets, university of chicago technical report 90-22.
\newblock {\em Computer Science Department}, 20, 1990.

\bibitem{fujii2015power}
Keisuke Fujii, Hirotada Kobayashi, Tomoyuki Morimae, Harumichi Nishimura, Shuhei Tamate, and Seiichiro Tani.
\newblock Power of quantum computation with few clean qubits.
\newblock {\em arXiv preprint arXiv:1509.07276}, 2015.

\bibitem{han1997threshold}
Yenjo Han, Lane~A Hemaspaandra, and Thomas Thierauf.
\newblock Threshold computation and cryptographic security.
\newblock {\em SIAM Journal on Computing}, 26(1):59--78, 1997.

\bibitem{morimae2017hardness}
Tomoyuki Morimae.
\newblock Hardness of classically sampling the one-clean-qubit model with constant total variation distance error.
\newblock {\em Physical Review A}, 96(4):040302, 2017.

\bibitem{o2018weakness}
Ryan O'Donnell and AC~Cem Say.
\newblock The weakness of ctc qubits and the power of approximate counting.
\newblock {\em ACM Transactions on Computation Theory (TOCT)}, 10(2):1--22, 2018.

\bibitem{stockmeyer1985approximation}
Larry Stockmeyer.
\newblock On approximation algorithms for\# p.
\newblock {\em SIAM Journal on Computing}, 14(4):849--861, 1985.

\bibitem{toda1991pp}
Seinosuke Toda.
\newblock Pp is as hard as the polynomial-time hierarchy.
\newblock {\em SIAM Journal on Computing}, 20(5):865--877, 1991.

\bibitem{valiant1985np}
Leslie~G Valiant and Vijay~V Vazirani.
\newblock Np is as easy as detecting unique solutions.
\newblock In {\em Proceedings of the seventeenth annual ACM symposium on Theory of computing}, pages 458--463, 1985.

\end{thebibliography}
\end{document}
